\section{Picard groups of rings}
\label{section-picard}

\noindent
We first define invertible modules as follows.

\begin{definition}
\label{definition-invertible}
Let $R$ be a ring. An $R$-module $M$ is {\it invertible} if the functor
$$
\text{Mod}_R \longrightarrow \text{Mod}_R,\quad
N \longmapsto M \otimes_R N
$$
is an equivalence of categories. An invertible $R$-module is said to be
{\it trivial} if it is isomorphic to $R$ as an $R$-module.
\end{definition}

\begin{lemma}
\label{lemma-invertible}
Let $R$ be a ring. Let $M$ be an $R$-module. Equivalent are
\begin{enumerate}
\item $M$ is finite locally free module of rank $1$,
\item $M$ is invertible, and
\item there exists an $R$-module $N$ such that $M \otimes_R N \cong R$.
\end{enumerate}
Moreover, in this case the module $N$ in (3) is isomorphic
to $\Hom_R(M, R)$.
\end{lemma}

\begin{proof}
Assume (1). Consider the module $N = \Hom_R(M, R)$ and the evaluation
map $M \otimes_R N = M \otimes_R \Hom_R(M, R) \to R$. If $f \in R$
such that $M_f \cong R_f$, then the evaluation map becomes an isomorphism
after localization at $f$ (details omitted). Thus we see the evaluation
map is an isomorphism by Algebra, Lemma \ref{algebra-lemma-cover}.
Thus (1) $\Rightarrow$ (3).

\medskip\noindent
Assume (3). Then the functor $K \mapsto K \otimes_R N$ is a quasi-inverse
to the functor $K \mapsto K \otimes_R M$. Thus (3) $\Rightarrow$ (2).
Conversely, if (2) holds, then $K \mapsto K \otimes_R M$ is essentially
surjective and we see that (3) holds.

\medskip\noindent
Assume the equivalent conditions (2) and (3) hold. Denote
$\psi : M \otimes_R N \to R$ the isomorphism from (3).
Choose an element $\xi = \sum_{i = 1, \ldots, n} x_i \otimes y_i$
such that $\psi(\xi) = 1$. Consider the isomorphisms
$$
M \to M \otimes_R M \otimes_R N \to M
$$
where the first arrow sends $x$ to $\sum x_i \otimes x \otimes y_i$
and the second arrow sends $x \otimes x' \otimes y$ to $\psi(x' \otimes y)x$.
We conclude that $x \mapsto \sum \psi(x \otimes y_i)x_i$ is
an automorphism of $M$. This automorphism factors as
$$
M \to R^{\oplus n} \to M
$$
where the first arrow is given by
$x \mapsto (\psi(x \otimes y_1), \ldots, \psi(x \otimes y_n))$
and the second arrow by $(a_1, \ldots, a_n) \mapsto \sum a_i x_i$.
In this way we conclude that $M$ is a direct summand of a finite free
$R$-module. This means that $M$ is finite locally free
(Algebra, Lemma \ref{algebra-lemma-finite-projective}).
Since the same is true for $N$ by symmetry and since
$M \otimes_R N \cong R$, we see that
$M$ and $N$ both have to have rank $1$.
\end{proof}

\noindent
The set of isomorphism classes of these
modules is often called the {\it class group} or {\it Picard group}
of $R$. The group structure is determined by assigning to
the isomorphism classes of the invertible modules $L$ and $L'$
the isomorphism class of $L \otimes_R L'$.
The inverse of an invertible module $L$ is the module
$$
L^{\otimes -1} = \Hom_R(L, R),
$$
because as seen in the proof of Lemma \ref{lemma-invertible}
the evaluation map $L \otimes_R L^{\otimes -1} \to R$ is an isomorphism.
Let us denote the Picard group of $R$ by $\Pic(R)$.

\begin{lemma}
\label{lemma-UFD-Pic-trivial}
Let $R$ be a UFD. Then $\Pic(R)$ is trivial.
\end{lemma}

\begin{proof}
Let $L$ be an invertible $R$-module. By Lemma \ref{lemma-invertible}
we see that $L$ is a finite locally free $R$-module. In particular
$L$ is torsion free and finite over $R$. Pick a nonzero element
$\varphi \in \Hom_R(L, R)$ of the dual invertible module.
Then $I = \varphi(L) \subset R$ is an ideal which is an invertible module.
Pick a nonzero $f \in I$ and let
$$
f = u p_1^{e_1} \ldots p_r^{e_r}
$$
be the factorization into prime elements with $p_i$ pairwise distinct.
Since $L$ is finite locally free there exist $a_i \in R$,
$a_i \not \in (p_i)$ such that $I_{a_i} = (g_i)$ for some $g_i \in R_{a_i}$.
Then $p_i$ is still a prime element of the UFD $R_{a_i}$ and
we can write $g_i = p_i^{c_i} g'_i$ for some $g'_i \in R_{a_i}$
not divisible by $p_i$. Since $f \in I_{a_i}$ we see that $e_i \geq c_i$.
We claim that $I$ is generated by $h = p_1^{c_1} \ldots p_r^{c_r}$ which
finishes the proof.

\medskip\noindent
To prove the claim it suffices to show that $I_a$ is generated by $h$
for any $a \in R$ such that $I_a$ is a principal ideal
(Algebra, Lemma \ref{algebra-lemma-cover}). Say $I_a = (g)$.
Let $J \subset \{1, \ldots, r\}$ be the set of $i$ such that
$p_i$ is a nonunit (and hence a prime element) in $R_a$. Because
$f \in I_a = (g)$ we find the prime factorization
$g = v \prod_{i \in J} p_j^{b_j}$
with $v$ a unit and $b_j \leq e_j$. For each $j \in J$ we have
$I_{aa_j} = g R_{aa_j} = g_j R_{aa_j}$, in other words
$g$ and $g_j$ map to associates in $R_{aa_j}$. By uniqueness
of factorization this implies that $b_j = c_j$ and the proof is complete.
\end{proof}







